\scanpage{96}\mspage{87}
Posons
\[
 \delta X=\left\{\begin{aligned}
 &\unsure{x|X}\ \text{si }\deleted{\unsure{x}}X=P_x\\
 &|\partial X|=\{a,b\}\ \text{si }X=S_{a,b}
 \end{aligned}\right.
\]
donc en tous cas
\[
 \delta X\subset L,\qquad \unsure{\operatorname{card}|\operatorname{Card}}\delta X\in\{1,2\}
\]

\deleted{Soient} Une condition \uline{nécessaire} pour $X \mathrel{\vcenter{\hbox{\tikz[x=.8em,y=.5em,line join=round]{\draw (0,1)--(1,.5)--(0,0) (1,1)--(0,.5)--(1,0);}}}} Y$ est
\[
\begin{aligned}
 X \mathrel{\vcenter{\hbox{\tikz[x=.8em,y=.5em,line join=round]{\draw (0,1)--(1,.5)--(0,0) (1,1)--(0,.5)--(1,0);}}}} Y\Longrightarrow\delta(\unsure{x|X})\cup\delta(Y)\in
  \underset{\inserted{\substack{\text{ens des parties}\\\text{finies }\deleted{\unsure{to}}\text{tot. ord de }L.}}}
            {\underline{\mathrm{Drap}(L)}}
\end{aligned}
\]

\deleted{\unsure{cette}} ce drapeau peut avoir cardinal égal à $1$. (si $X=Y\in M_0$), $2,3$, ou $4$ (si $X,Y\in M_1$ et $X\parallel Y$). Le cas de cardinal $2$ est celui de $P_{\unsure{a|x}} \mathrel{\vcenter{\hbox{\tikz[x=.8em,y=.5em,line join=round]{\draw (0,1)--(1,.5)--(0,0) (1,1)--(0,.5)--(1,0);}}}} P_y$ avec $x\ne y$, ou de $P_x \mathrel{\vcenter{\hbox{\tikz[x=.8em,y=.5em,line join=round]{\draw (0,1)--(1,.5)--(0,0) (1,1)--(0,.5)--(1,0);}}}} S_{a,b}$ avec $P_x\triangleleft S_{ab}$ i.e. $x\in\{a,b\}$, ou $\deleted{\unsure{S}}X=Y\in\deleted{\unsure{S}}M_1$\,; le cas de cardinal $3$ est $X\in M_0$, $Y\in M_1$ ou inversement, avec $X\parallel Y$ (soit $P_x,S_{a,b}$ avec $x<a$ ou $\unsure{x|a}>b$), ou celui de $X,Y\in M_1$, placés comme $\deleted{\unsure{S_{a,}}}\ S_{a,b},S_{b,c}$ avec $a<b<c$ -- (segments \uline{contigus}). Soit tabl. possibilités
\diagram{\(\displaystyle
 \left\{\begin{array}{@{}l@{}}
 \text{lieux égaux ou }\deleted{\text{\unsure{intérieurement}}}\ \text{disjoints}\\
 \text{lieu }\deleted{\text{adjacent}}\ \text{incident à un segment, ou disjoint du}\\
 \hspace{17em}\text{segment}\\
 \text{segments disjoints, segments adjacents}
 \end{array}\right.
\)}

Mais de façon générale, si on a $T\in\mathrm{Drap}(L)$,
\[
 T=\{t_1,t_2,\ldots,t_{\unsure{n|m}}\}\quad\text{avec}\quad t_1<t_2<\ldots<t_n,
\]
on \unsure{définit|définira} les \uline{segments} \uline{interstitiels} \inserted{de $T$} comme segments $\deleted{\unsure{P}}S_{t_1,t_2},S_{t_2,t_3},\ldots S_{t_{n-1},t_n}$. \deleted{Ceci}\inserted{les} \deleted{dit, on trouve ceci.} $P_{t_i}$ sont appelés les \uline{lieux} interstitiels. Ceci dit, on trouve ceci\,:
% slanted
\margin{\uline{NB} Si $\operatorname{card}T=1$, il n'y a pas \unsure{de segment interstitiel|de segments interstitiels}, \unsure{seult} \uline{un} lieu interstitiel}
% archive sheet number: 96

